\chapter{Complexity — Random Walk Mixing Pad}
%%% generated by the blueprint pipeline from TCSlib.Complexity.Randomized.RandomWalkMixingPad
\section{Overview}

Declarations of the Lean module \texttt{TCSlib.Complexity.Randomized.RandomWalkMixingPad}.

\section{Declarations}

\begin{theorem}[Bounds on the hitting probability]
\lean{Randomized.hitProb_mk_mem_Icc}
\label{Randomized.hitProb_mk_mem_Icc}
\leanok
\uses{Randomized.avgOut, Randomized.hitProb, Randomized.outNbrs}
\difficulty{5}
Let $n$ be a natural number, let $r$ be a decidable binary relation on $\text{Fin}(n)$,
and form the directed graph whose edges are given by $r$. For every target vertex $t$,
time $k$, and starting vertex $s$, the probability of hitting $t$ after $k$ steps is
between $0$ and $1$, inclusive.
\end{theorem}

\begin{theorem}[Zero hitting probability from an isolated vertex]
\lean{Randomized.hitProb_mk_eq_zero}
\label{Randomized.hitProb_mk_eq_zero}
\leanok
\uses{Randomized.avgOut, Randomized.hitProb, Randomized.outNbrs}
\difficulty{4}
Let $n$ be a natural number and let $r$ be a decidable relation on $\text{Fin}(n)$. For
distinct vertices $s,t\in\text{Fin}(n)$, if $s$ has no outgoing neighbors in the digraph
defined by $r$, then for every natural number $k$, the $k$-step probability of hitting
$t$ from $s$ is zero.
\end{theorem}

\begin{theorem}[Hitting probability recurrence]
\lean{Randomized.hitProb_mk_succ}
\label{Randomized.hitProb_mk_succ}
\leanok
\uses{Randomized.avgOut, Randomized.hitProb, Randomized.outNbrs}
\difficulty{1}
Let $n$ be a natural number, let $r$ be a decidable binary relation on $\text{Fin}(n)$,
and form the directed graph whose edges are given by $r$. If $s$ and $t$ are distinct
vertices and $s$ has at least one outgoing neighbor, then for every $k\in\mathbb{N}$ the
probability of hitting $t$ from $s$ after $k+1$ steps is the average of the
probabilities after $k$ steps over all outgoing neighbors of $s$.
\end{theorem}

\begin{theorem}[Sum over padded neighbors]
\lean{Randomized.sum_padNbr}
\label{Randomized.sum_padNbr}
\leanok
\uses{Randomized.padNbr}
\difficulty{5}
Let $n,D$ be natural numbers, let $r$ be a decidable relation on the set of $n$ indices,
and suppose $n\le D$. For any rational-valued function $f$ on those indices and any
index $s$, the sum of $f$ over the $D$ padded neighbors of $s$ equals the sum of $f$
over the indices related to $s$, plus $D$ minus the number of such related indices,
multiplied by $f(s)$.
\end{theorem}

\begin{definition}[Binomial tail recursion]
\lean{Randomized.binTail}
\label{Randomized.binTail}
\leanok
For a relation $r$ on $\\{0,\\ldots,n-1\\}$ and a natural number $D$, define $B_D(m,j)$
recursively by $B_D(m,0)=0$, $B_D(0,j+1)=1$, and
$B_D(m+1,j+1)=B_D(m,j)/D+(1-1/D)B_D(m,j+1)$.
\end{definition}

\begin{theorem}[Reciprocal bound]
\lean{Randomized.one_div_mem}
\label{Randomized.one_div_mem}
\leanok
\difficulty{3}
For every natural number $D$ with $1 \le D$, the rational number $1/D$ is positive and
at most $1$.
\end{theorem}

\begin{theorem}[Bounds for the binomial tail recurrence]
\lean{Randomized.binTail_mem_Icc}
\label{Randomized.binTail_mem_Icc}
\leanok
\uses{Randomized.binTail, Randomized.one_div_mem}
\difficulty{4}
Let $D$ be a natural number with $1\le D$, and define $T_D(m,0)=0$, $T_D(0,j+1)=1$, and
$T_D(m+1,j+1)=T_D(m,j)/D+(1-1/D)T_D(m,j+1)$. For every pair of natural numbers $m,j$,
one has $0\le T_D(m,j)\le 1$.
\end{theorem}

\begin{theorem}[Monotonicity in the tail threshold]
\lean{Randomized.binTail_mono}
\label{Randomized.binTail_mono}
\leanok
\uses{Randomized.binTail, Randomized.binTail_mem_Icc, Randomized.one_div_mem}
\difficulty{4}
Let $D$ be a natural number with $1\le D$, and define $T_D(m,0)=0$, $T_D(0,j+1)=1$, and
$T_D(m+1,j+1)=T_D(m,j)/D+(1-1/D)T_D(m,j+1)$. For all natural numbers $m,j$, the sequence
is nondecreasing in its second argument: $T_D(m,j)\le T_D(m,j+1)$.
\end{theorem}

\begin{theorem}[Monotonicity in the number of trials]
\lean{Randomized.binTail_succ_le}
\label{Randomized.binTail_succ_le}
\leanok
\uses{Randomized.binTail, Randomized.binTail_mono, Randomized.one_div_mem}
\difficulty{3}
Let $D$ be a natural number with $1\le D$, and define $T_D(m,0)=0$, $T_D(0,j+1)=1$, and
$T_D(m+1,j+1)=T_D(m,j)/D+(1-1/D)T_D(m,j+1)$. For all natural numbers $m,j$, increasing
the first argument by one does not increase the value: $T_D(m+1,j)\le T_D(m,j)$.
\end{theorem}

\begin{theorem}[Bound for sums of binomial tails]
\lean{Randomized.sum_binTail_le}
\label{Randomized.sum_binTail_le}
\leanok
\uses{Randomized.binTail, Randomized.one_div_mem}
\difficulty{4}
Let $D$ be a positive natural number. For all natural numbers $M$ and $j$, the sum of
the first $M$ values of the recursively defined binomial tail with parameters $D$ and
$j$ is at most $jD$.
\end{theorem}

\begin{theorem}[Bound for the binomial-tail recurrence]
\lean{Randomized.succ_mul_binTail_le}
\label{Randomized.succ_mul_binTail_le}
\leanok
\uses{Randomized.binTail, Randomized.binTail_succ_le, Randomized.sum_binTail_le}
\difficulty{3}
Let $D$ be a natural number with $1\le D$, and define $B_D$ by $B_D(m,0)=0$,
$B_D(0,j+1)=1$, and $B_D(m+1,j+1)=B_D(m,j)/D+(1-1/D)B_D(m,j+1)$. For all natural numbers
$L$ and $j$, $(L+1)B_D(L,j)\le jD$.
\end{theorem}

\begin{theorem}[Binomial-tail lower bound for hitting probability]
\lean{Randomized.sub_binTail_le_hitProb}
\label{Randomized.sub_binTail_le_hitProb}
\leanok
\uses{Randomized.RegularMultigraph.hitProb, Randomized.RegularMultigraph.hitProb_nonneg, Randomized.RegularMultigraph.hitProb_succ, Randomized.RegularMultigraph.hitProb_zero, Randomized.binTail, Randomized.binTail_mem_Icc, Randomized.binTail_mono, Randomized.hitProb, Randomized.hitProb_mk_eq_zero, Randomized.hitProb_mk_mem_Icc, Randomized.hitProb_mk_succ, Randomized.outNbrs, Randomized.padGraph, Randomized.padNbr, Randomized.sum_padNbr}
\difficulty{7}
Let $V$ be a finite set, let $G$ be a directed graph on $V$, and define the
out-neighbors of a vertex by its directed edges. For a symmetric relation $r$ on
$\{0,\ldots,n-1\}$, let $B_D(m,j)$ be given by $B_D(m,0)=0$, $B_D(0,j+1)=1$, and
$B_D(m+1,j+1)=B_D(m,j)/D+(1-1/D)B_D(m,j+1)$; let the padded graph use $D$ edge slots per
vertex, sending slots beyond the first $n$ to the vertex itself. For every $m,j$ and
starting vertex $s$, the hitting probability of a target $t$ in $G$ after $j$
independent uniform out-neighbor steps, minus $B_D(m,j)$, is at most the hitting
probability of $t$ after $m$ such steps in the padded graph, provided $n\le D$.
\end{theorem}

\begin{theorem}[Hit probability after bounded degree padding]
\lean{Randomized.half_le_hitProb_padGraph}
\label{Randomized.half_le_hitProb_padGraph}
\leanok
\uses{Randomized.RegularMultigraph.hitProb, Randomized.binTail, Randomized.binTail_mem_Icc, Randomized.hitProb, Randomized.padGraph, Randomized.sub_binTail_le_hitProb, Randomized.succ_mul_binTail_le}
\difficulty{3}
Let $n,D$ be nonnegative integers and let $r$ be a symmetric relation on $\text{Fin}(n)$,
with $n\le D\le 4n^2$. For vertices $s,t$ and an integer $k$, if the probability that a
length-$k$ walk in the directed graph defined by $r$, starting at $s$, visits $t$ is at
least $9/10$, then the probability that a length-$10n^2k$ walk in the padded regular
multigraph of degree $D$, starting at $s$, visits $t$ is at least $1/2$.
\end{theorem}

\begin{theorem}[Hit probability after degree padding]
\lean{Randomized.half_le_hitProb_padDegree}
\label{Randomized.half_le_hitProb_padDegree}
\leanok
\uses{Randomized.RegularMultigraph.hitProb, Randomized.half_le_hitProb_padGraph, Randomized.hitProb, Randomized.le_padDegree, Randomized.padDegree_le, Randomized.padGraph}
\difficulty{3}
Let $n$ be a nonnegative integer and let $r$ be a symmetric relation on $\text{Fin}(n)$.
For vertices $s,t$ and an integer $k$, if the probability that a length-$k$ walk in the
directed graph defined by $r$, starting at $s$, visits $t$ is at least $9/10$, then the
corresponding probability in the padded regular multigraph of degree
$	exttt{le\_padDegree}(n)$, after $10n^2k$ steps, is at least $1/2$.
\end{theorem}
